% \subsubsection{Mailbox Types}
% \label{sec:mailbox-types-intro}

Where session types reason about channels between processes \cite{hondaTypesDyadicInteraction1993,wadlerPropositionsSessions2014,kobayashiTypeSystemLockFree2002,huttelFoundationsSessionTypes2016}, \emph{mailbox types} \cite{deliguoroMailboxTypesUnordered2018} reason about mailboxes.
% Mailboxes are mostly present in actor-based programming languages such as Erlang and Elixir.
% There, a process is able to retrieve messages from its own mailbox and send messages to other mailboxes.
% The main difference to channels is the order in which a process retrieves messages.
% A process may delay processing a message before some other message arrives, and in general is not bound by the order of message arrivals.

Similar to how the $\pi$-calculus \cite{milnerCalculusMobileProcesses1992} models process communication via channels, the \emph{mailbox calculus},
introduced by de’Liguoro and Padovani as an extension of the asynchronous $\pi$-calculus \cite{amadioBisimulationsAsynchronousPcalculus1998},
models process communication via message passing and mailboxes.
Mailbox names are treated as first-class objects that can be sent to other mailboxes in messages.

\emph{Mailbox types} are introduced as a way of reasoning about the content of mailboxes.
In the context of the mailbox calculus, mailbox types eliminate protocol violations and deadlocks, e.g.\ by keeping track of the dependencies
between mailbox names using a dependency graph and checking for acyclicity.

% TODO: Keep example or remove it?
% As an example of how terms in the mailbox calculus look like, the authors provide a definition for a \emph{future variable}, representing a delayed computation \cite{deliguoroMailboxTypesUnordered2018,fowlerSpecialDeliveryProgramming2023}.
% % Figure \ref{fig:future-protocol} shows the general idea of the protocol, using two clients and one server process.
% One process sends a $\msg{Put}$ message with some value $x$ to the server.
% Then, other processes may retrieve this message by sending a $\msg{Get}$ message.
% The server should only be receiving a single $\msg{Put}$ message at the beginning, while being able to respond to an arbitrary amount of $\msg{Get}$ messages afterwards.
%
% % \begin{figure}[t]
% %   \centering
% % \begin{sequencediagram}
% %   \tikzstyle{dotted}=[line width=1pt]
% %   \newinst{c1}{Client 1}
% %   \newinst[1]{s}{Server}
% %   \newinst[1]{c2}{Client 2}
% %   \mess{c1}{$\msg{Put}[x]$}{s}
% %   \postlevel
% %   \mess{c1}{$\msg{Get}[C_1]$}{s}
% %   \mess{s}{$\msg{Reply}[x]$}{c1}
% %   \node[anchor=east] (t0) at (mess from) {};
% %   \node[anchor=west] (t1) at (mess to) {};
% %   \coordinate (leftC) at ($(t0)!0.5!(t1)$);
% %   \node[yshift=-0.4cm] (leftDots) at (leftC) {$\vdots$};
% %   \postlevel
% %   \mess{c2}{$\msg{Get}[C_2]$}{s}
% %   \mess{s}{$\msg{Reply}[x]$}{c2}
% % \end{sequencediagram}
% %   \caption{Protocol diagram of the future example}
% %   \label{fig:future-protocol}
% % \end{figure}
% %
% % The idea behind this protocol can be described in terms of the mailbox calculus, where $\mathit{self}$ is a mailbox name.
%
% \vspace{-0.4cm}
% \small
% \begin{align*}
%   \mathsf{emptyFuture}(\mathit{self}) &\triangleq \mathit{self}\mathsf{?}\msg{Put}(x).\mathsf{fullFuture}[\mathit{self}, x]\\
%   \mathsf{fullFuture}(\mathit{self}, x) &\triangleq \keyword{free } \mathit{self}.\keyword{done}\\
%   &+ \mathit{self}\mathsf{?}\msg{Get}(sender).(sender\mathsf{!}\msg{Reply}[x] \parallel \mathsf{fullFuture}[\mathit{self}, x])\\
%   &+ \mathit{self}\mathsf{?}\msg{Put}.\keyword{fail } \mathit{self}
% \end{align*}
% \begin{align*}
%   &(\nu \mathit{future})(\mathsf{emptyFuture}(\mathit{future}) \parallel future\mathsf{!}\msg{Put}[5] \parallel\\
%   &(\nu \mathit{\mathit{self}})(future\mathsf{!}\msg{Get}[\mathit{self}] \parallel (\mathit{self}\mathsf{?}\msg{Reply}(x).\keyword{free} \ \mathit{self}.\mathsf{print}(\mathsf{intToString}(x)))))
% \end{align*}
% \normalsize
%
% The process defined by $\mathsf{emptyFuture}$ will wait for message $\msg{Put}$ with payload $x$ to be present in mailbox $\mathit{self}$ and then continues as specified by $\mathsf{fullFuture}$.
% For $\mathsf{fullFuture}$ there are three possible actions: free mailbox $\mathit{self}$ if it is contains no messages and be done, fail if it retrieves another $\msg{Put}$ message,
% or retrieving a $\msg{Get}$ message, it will send a $\msg{Reply}$ message with $x$ to the sender and continues recursively.
% This behaviour captures the notion of the value only being computed once, but retrieved arbitrarily often.
% The last term creates a new mailbox name $\mathit{future}$ and calls $\mathsf{emptyFuture}$ with it as an argument.
% In parallel, the message $\msg{Put}$ is sent to $\mathit{future}$, while also creating new mailbox name $\mathit{self}$.
% After $\mathit{self}$ is created, a $\msg{Get}$ message is sent to $\mathit{future}$ with $\mathit{self}$ as mailbox to respond to.
% Finally, when $\mathit{self}$ receives a $\msg{Reply}$ message, the process frees $\mathit{self}$, printing response $x$.

% TODO: Keep example or remove it?
% By defining the appropriate types, the above example of a future variable successfully type checks.
% The type of $\mathsf{emptyFuture}$ can be defined as $\mbtin{(\msg{Put}[\mathsf{Int}] \mbcomp \mbstar{\msg{Get}}[\mbtout{\msg{Reply}[\mathsf{Int}]}])}$.
% The mailbox can contain a $\msg{Put}$ message with an integer payload and 
% an arbitrary amount of $\msg{Get}$ messages with a payload of a reference to a mailbox that expects a $\msg{Reply}$ message.
% The type of $\mathsf{fullFuture}$ can be defined as $\mbtin{(\mbstar{\msg{Get}}[\mbtout{\msg{Reply}[\mathsf{Int}]}])}$, indicating that the mailbox
% does not longer contain a $\msg{Put}$ message and only consists of zero or more $\msg{Get}$ messages.
% Each of these $\msg{Get}$ messages is required to be paired with a response by sending a $\msg{Reply}$ message.

